\section{9.3 정규화 기법과 학습률 스케줄링}
\begin{frame}{과적합과 신경망 전용 정규화}
  \begin{itemize}
    \item \kb{과적합}: 학습 데이터는 완벽히 맞히는데 새 데이터에서 성능
      저하 — Week 4.1 편향-분산 관점에서, 파라미터가 매우 많은 신경망은
      \textbf{극단적으로 유연(분산 큼)}해 특히 취약
    \item Week 6의 L1/L2 정규화는 그대로 쓸 수 있으나, 신경망은
      \textbf{구조 자체를 이용}한 정규화도 함께 사용:
  \end{itemize}
  \vskip0.3em
  \begin{columns}[T]
    \begin{column}{0.5\textwidth}
      \kb{Dropout}
      \begin{itemize}
        \item 학습 시 매 스텝마다 일부 뉴런을 확률 $p$로
          \textbf{무작위로 꺼버림}
        \item 뉴런 간 \textbf{공적응(co-adapt)}을 억제 — 특정
          뉴런에만 의존하는 대신 \textbf{견고한(robust)} 특징 학습
      \end{itemize}
    \end{column}
    \begin{column}{0.5\textwidth}
      \kb{Batch Normalization (BN)}
      \begin{itemize}
        \item 각 층의 입력을 \textbf{미니배치} 단위로 평균 0,
          분산 1로 정규화, 이후 학습 가능한 $\gamma, \beta$를
          다시 곱·추가
        \item \kb{internal covariate shift} 감소 $\to$ 학습이
          \textbf{더 빠르고 안정적}
      \end{itemize}
    \end{column}
  \end{columns}
  \vskip0.3em
  {\footnotesize ``정규화(regularization)''라는 말은 릿지 회귀
  (Tikhonov, 1963)에서 왔고, 2010년대 dropout·BN 등장 이후
  ``손실을 바꾸는 L1/L2'' + ``\textbf{계산 구조 자체를 바꾸는}
  구조적 정규화''라는 새 축이 생김.}
\end{frame}
